% Generated draft: requires mathematical and source review.
\documentclass[11pt]{article}
\usepackage[margin=25mm]{geometry}
\usepackage{fontspec}
\setmainfont{Latin Modern Roman}
\usepackage{amsmath,amssymb,mathtools,unicode-math}
\setmathfont{Latin Modern Math}
\usepackage{enumitem,needspace,xurl,hyperref}
\setlength{\parindent}{0pt}
\setlength{\parskip}{5pt}
\setlength{\emergencystretch}{3em}
\newcommand{\N}{\mathbb{N}}
\newcommand{\Z}{\mathbb{Z}}
\newcommand{\Q}{\mathbb{Q}}
\newcommand{\R}{\mathbb{R}}
\newcommand{\C}{\mathbb{C}}
\newcommand{\F}{\mathbb{F}}
\newcommand{\A}{\mathbb{A}}
\newcommand{\PP}{\mathbb{P}}
\newcommand{\E}{\mathbb{E}}
\newcommand{\eps}{\varepsilon}
\newcommand{\dd}{\,\mathrm{d}}
\begin{document}
\section*{UnsolvedMath ID: 30005677}
\textbf{Source problem code:} catalogue entry
\subsection{Definitions and mathematical statement}
Let \(T(n)\) denote the chromatic telescope intended by the source, and let \(L_{T(n)}\) denote Bousfield localization of spectra with respect to \(T(n)\). A spectrum \(X\) is \(T(n)\)-local if its localization map \(X\to L_{T(n)}X\) is an equivalence. The supplied record does not specify the prime, the permitted heights \(n\), or the telescope data; these remain parameters requiring confirmation from the source.

Write \(\mathcal L_{T(n)}\) for the intended \(\infty\)-category of \(T(n)\)-local spectral Lie algebras, \(\mathcal H_{T(n)}\) for the intended \(\infty\)-category of \(T(n)\)-local spectral Hopf algebras, and
\[
U\colon \mathcal L_{T(n)}\longrightarrow \mathcal H_{T(n)}
\]
for the universal enveloping algebra functor equipped with the Hopf structure intended by the source. Here \(\operatorname{Map}_{\mathcal C}(A,B)\) denotes the mapping space between objects \(A,B\) of an \(\infty\)-category \(\mathcal C\).

\textbf{Conjectural statement.} For every pair \(\mathfrak g,\mathfrak h\in\mathcal L_{T(n)}\), the map induced by \(U\),
\[
\operatorname{Map}_{\mathcal L_{T(n)}}(\mathfrak g,\mathfrak h)
\longrightarrow
\operatorname{Map}_{\mathcal H_{T(n)}}(U\mathfrak g,U\mathfrak h),
\]
is an equivalence of spaces. Thus \(U\) is fully faithful in the \(\infty\)-categorical sense, including all higher mapping information.

\textbf{Unresolved formulation qualifications.} This expresses the full-faithfulness assertion in the supplied statement, conditional on the intended categories and functor. The record does not identify the spectral Lie operad, specify the Hopf algebra structure and its morphisms, or state whether cocommutativity, an antipode, completion, or additional restrictions on objects are required. Consequently, these symbols do not yet determine a unique fully specified conjecture. The supplied original excerpt calls the intended assertion ``Conjecture A'' and a parallel of the Milnor--Moore theorem, but does not reproduce its mathematical content.

\textbf{Supplied status qualification.} The record labels the problem partially solved. Its dated assessment states that full faithfulness holds on abelian \(T(n)\)-local spectral Lie algebras, including infinite-loop cases, while the assertion for all \(T(n)\)-local spectral Lie algebras remains unproved. These are supplied curation claims, not independently verified results here; the meanings and hypotheses of the stated special cases require source confirmation.
\subsection{Short English statement}
The conjecture says that universal enveloping algebras, with their intended spectral Hopf structure, retain the entire mapping space between any two \(T(n)\)-local spectral Lie algebras. The source record leaves the precise categories and chromatic parameters unspecified. Its supplied partial-solution assessment covers abelian and infinite-loop cases, subject to independent verification.
\subsection{Sources}
\begin{itemize}
\item Local UnsolvedMath registry, permanent ID \(30005677\), title: Fully Faithful Enveloping Algebras for Spectral Lie Algebras; source problem code \texttt{catalogue entry}. Collection: Oberwolfach Reports Open Problems. No individual proposer is identified in the supplied record; the proposed year is recorded as 2023.
\item Supplied source citation: \emph{Homotopy Theory (2024)}, Oberwolfach Reports, pp.\ 1893--1952, DOI \texttt{10.4171/owr/2023/34}. \url{https://doi.org/10.4171/owr/2023/34}. The record describes Oberwolfach Reports as workshop reports containing open-problem discussions and identifies the intended assertion as Conjecture A, a parallel of the Milnor--Moore theorem. The cited title year is 2024, while the DOI and proposed-year field indicate 2023; their relationship is not explained in the supplied record.
\item Supplied literature source labelled ``Antolin-Camarena-Brantner-Heuts'': \url{https://arxiv.org/abs/2501.03116}. No title, theorem number, or precise supporting passage is supplied.
\item Literature assessment checked at \texttt{2026-08-22}: supplied dated curation, not independently verified. The record reports full faithfulness on abelian \(T(n)\)-local spectral Lie algebras, including infinite-loop cases, and reports that the general conjecture is not proved. Its listed literature sources are the Oberwolfach report and the Antolin-Camarena-Brantner-Heuts link above. The background explicitly requests independent verification of time-sensitive status and cited claims.
\item The registry labels the statement \texttt{corrected\_verified} and says it was reconstructed and checked against the official Oberwolfach report. This is supplied verification metadata, not a source check performed for this draft.
\end{itemize}
\end{document}
